\documentclass[11pt]{article}
\usepackage[margin=18mm,headheight=22pt]{geometry}
\usepackage{amsmath}
\usepackage{fontspec}
\usepackage{unicode-math}
\setmainfont{NewCM10}[Extension=.otf,UprightFont=*-Book,BoldFont=*-Bold,ItalicFont=*-BookItalic,BoldItalicFont=*-BoldItalic]
\setsansfont{NewCMSans10}[Extension=.otf,UprightFont=*-Book,BoldFont=*-Bold,ItalicFont=*-BookOblique,BoldItalicFont=*-BoldOblique]
\setmathfont{NewCMMath-Book.otf}
\usepackage{xcolor}
\definecolor{FAUblue}{HTML}{04316A}
\definecolor{Paperblue}{HTML}{EEF2F7}
\definecolor{EEIgold}{HTML}{CC9C42}
\usepackage{etoolbox,enumitem,booktabs,tabularx,array,needspace,microtype}
\usepackage[most]{tcolorbox}
\usepackage{titlesec,fancyhdr}
\usepackage[hidelinks]{hyperref}
\usepackage{draftwatermark}
\SetWatermarkText{Made by Thanate Amornjiraporn (Focus)}
\SetWatermarkColor{black!7}\SetWatermarkFontSize{30pt}\SetWatermarkAngle{45}
\setlength{\parindent}{0pt}
\setlength{\parskip}{5pt}
\setlist{itemsep=2pt,topsep=3pt,leftmargin=1.6em,beginpenalty=10000,endpenalty=10000}
\raggedbottom
\clubpenalty=10000 \widowpenalty=10000 \displaywidowpenalty=10000 \interlinepenalty=10000
\setcounter{secnumdepth}{0}
\setcounter{tocdepth}{1}
\pagestyle{fancy}\fancyhf{}
\fancyhead[L]{\small\sffamily\color{FAUblue} From Fundamentals to AI Apocalypse · MEI · Exercise Sheet 4}
\fancyhead[R]{\small\sffamily Exercise 4: Worked solutions}
\fancyfoot[C]{\small\thepage}
\titleformat{\section}{\Large\bfseries\sffamily\color{FAUblue}}{}{0pt}{}
\titleformat{\subsection}{\large\bfseries\sffamily\color{FAUblue}}{}{0pt}{}
\titlespacing*{\section}{0pt}{16pt}{6pt}
\preto\section{\Needspace{8\baselineskip}}
\preto\subsection{\Needspace{6\baselineskip}}
\newtcolorbox{equationcard}[1]{colback=Paperblue,colframe=Paperblue,boxrule=0pt,arc=0pt,left=10pt,right=10pt,top=6pt,bottom=6pt,title={#1},fonttitle=\small\bfseries\sffamily,coltitle=FAUblue,before skip=8pt,after skip=8pt}
\newtcolorbox{workedbox}[1]{colback=white,colframe=EEIgold,boxrule=0pt,leftrule=3pt,arc=0pt,left=10pt,top=4pt,bottom=4pt,title={#1},fonttitle=\small\bfseries\sffamily,coltitle=FAUblue,attach title to upper={\par},before skip=8pt,after skip=8pt}
\newtcolorbox{recallbox}[1]{colback=white,colframe=FAUblue,boxrule=0pt,leftrule=3pt,arc=0pt,left=10pt,top=4pt,bottom=4pt,title={#1},fonttitle=\small\bfseries\sffamily,coltitle=FAUblue,attach title to upper={\par},before skip=8pt,after skip=8pt}
\usepackage{environ}
\newsavebox\measurebox
% Typeset the section once invisibly to measure it; start a new page if it will not fit in what is left.
\makeatletter
\newif\ifkeep@measuring
\NewEnviron{keepsection}{%
  \ifkeep@measuring \BODY \else
    \edef\keep@sec{\the\value{section}}\edef\keep@sub{\the\value{subsection}}%
    \setbox\measurebox\vbox{\hsize\linewidth\keep@measuringtrue\def\Needspace##1{}\BODY}%
    \setcounter{section}{\keep@sec}\setcounter{subsection}{\keep@sub}%
    \dimen2=\ht\measurebox \advance\dimen2\dp\measurebox
    \dimen@=\textheight
    \ifdim\pagegoal<\maxdimen
      \dimen@=\pagegoal \advance\dimen@-\pagetotal
      \ifdim\dimen2>\dimen@ \clearpage \dimen@=\textheight \fi
    \fi
    % if it fits, nothing inside may force an early break
    \ifdim\dimen2>\dimen@ \BODY \else \begingroup\def\Needspace##1{}\BODY\endgroup \fi
  \fi}
\makeatother
\newtcolorbox{notebox}{colback=Paperblue,colframe=Paperblue,boxrule=0pt,arc=0pt,left=8pt,right=8pt,top=4pt,bottom=4pt,before skip=6pt,after skip=6pt}
\newtcolorbox{plainbox}{colback=white,colframe=FAUblue,boxrule=0pt,leftrule=3pt,arc=0pt,left=10pt,right=8pt,top=4pt,bottom=4pt,before skip=6pt,after skip=8pt}
\usetikzlibrary{arrows.meta}
\title{\sffamily\bfseries\color{FAUblue}Exercise 4: Worked solutions}
\author{Mathematics for Engineers I \ensuremath{\cdot} Module 62765 \ensuremath{\cdot} Exercise Sheet 4\\[2pt]\small From Fundamentals to AI Apocalypse · personal study guide}
\date{}
\begin{document}
\maketitle
\begingroup\footnotesize\setlength{\parskip}{0pt}\tableofcontents\endgroup
\clearpage

\begin{keepsection}
\Needspace{7\baselineskip}
\section{Exercise 4.1: Mathematical induction}

\Needspace{4\baselineskip}
Induction starts from the first valid index and shows that each true case forces the next.

\textbf{Connected topic:} Induction (§2.5).

\begin{keepsection}
\Needspace{6\baselineskip}
\subsection{4.1a: Prove the finite geometric sum}

\Needspace{3\baselineskip}
Try first: Check n=0, then multiply the proposed sum by x and subtract, keeping x\ensuremath{\neq}1 visible.

\nopagebreak[4]{\small\itshape Answers: see “Exercise 4.1: Mathematical induction” under Answers at the end.}
\end{keepsection}

\begin{keepsection}
\Needspace{7\baselineskip}
\subsection{4.1b: Prove the reciprocal-square-root bound}

\Needspace{4\baselineskip}
Try first: Try induction from n=2: compare the new term with the change in \ensuremath{\sqrt{\mathrm{n}}}, or use a direct termwise bound.

\nopagebreak[4]{\small\itshape Answers: see “Exercise 4.1: Mathematical induction” under Answers at the end.}
\end{keepsection}
\end{keepsection}

\begin{keepsection}
\Needspace{9\baselineskip}
\section{Exercise 4.2: Maps}

\Needspace{6\baselineskip}
Composition applies f first and g second. For finite-set counts, choose an image for each input while respecting the required property.

\textbf{Connected topic:} Injectivity (§3.2) and finite counting (§2.4); Map types (§3.2) and counting and binomial coefficients (§2.4).

\begin{keepsection}
\Needspace{7\baselineskip}
\subsection{4.2a: Is a composition of injections injective?}

\Needspace{4\baselineskip}
Try first: Assume two composite outputs are equal and unwind the definition of injectivity in reverse order.

\nopagebreak[4]{\small\itshape Answers: see “Exercise 4.2: Maps” under Answers at the end.}
\end{keepsection}

\begin{keepsection}
\Needspace{6\baselineskip}
\subsection{4.2b(i): Count all maps M\ensuremath{\to}N}

\Needspace{3\baselineskip}
Try first: Each of the m elements of M independently chooses one of n target elements.

\nopagebreak[4]{\small\itshape Answers: see “Exercise 4.2: Maps” under Answers at the end.}
\end{keepsection}

\begin{keepsection}
\Needspace{7\baselineskip}
\subsection{4.2b(ii): Count bijections M\ensuremath{\to}N}

\Needspace{4\baselineskip}
Try first: A bijection requires equal cardinalities; when they match, count the possible orderings.

\nopagebreak[4]{\small\itshape Answers: see “Exercise 4.2: Maps” under Answers at the end.}
\end{keepsection}

\begin{keepsection}
\Needspace{7\baselineskip}
\subsection{4.2b(iii): Count injective maps M\ensuremath{\to}N}

\Needspace{4\baselineskip}
Try first: An injection assigns distinct targets to inputs, so there are fewer choices at each step.

\nopagebreak[4]{\small\itshape Answers: see “Exercise 4.2: Maps” under Answers at the end.}
\end{keepsection}

\begin{keepsection}
\Needspace{7\baselineskip}
\subsection{4.2b(iv): Count surjective maps M\ensuremath{\to}N}

\Needspace{4\baselineskip}
Try first: Exclude maps whose image misses at least one target, correcting for maps that miss several targets.

\nopagebreak[4]{\small\itshape Answers: see “Exercise 4.2: Maps” under Answers at the end.}
\end{keepsection}
\end{keepsection}

\begin{keepsection}
\Needspace{7\baselineskip}
\section{Exercise 4.3: Elementary functions}

\Needspace{4\baselineskip}
Addition formulas for sine and cosine imply the circle identity and double-angle relation.

\textbf{Connected topic:} Elementary functions and inverse functions (§3.5).

\begin{keepsection}
\Needspace{6\baselineskip}
\subsection{4.3a: Prove the circle equation}

\Needspace{3\baselineskip}
Try first: Expand cos(x\ensuremath{-}x) using the cosine addition theorem.

\nopagebreak[4]{\small\itshape Answers: see “Exercise 4.3: Elementary functions” under Answers at the end.}
\end{keepsection}

\begin{keepsection}
\Needspace{6\baselineskip}
\subsection{4.3b: Prove the cosine double-angle identity}

\Needspace{3\baselineskip}
Try first: Use cos(2x)=cos(x+x), then replace cos\textsuperscript{2}x by 1\ensuremath{-}sin\textsuperscript{2}x.

\nopagebreak[4]{\small\itshape Answers: see “Exercise 4.3: Elementary functions” under Answers at the end.}
\end{keepsection}

\begin{keepsection}
\Needspace{7\baselineskip}
\subsection{4.3c: Evaluate cos(arcsin x)}

\Needspace{4\baselineskip}
Try first: Let θ=arcsin x, so sin θ=x and θ lies in the principal interval where cosine is non-negative.

\nopagebreak[4]{\small\itshape Answers: see “Exercise 4.3: Elementary functions” under Answers at the end.}
\end{keepsection}

\begin{keepsection}
\Needspace{7\baselineskip}
\subsection{4.3d: Find cos(π/5) exactly}

\Needspace{4\baselineskip}
Try first: Use π/5=2(π/10) in the double-angle identity, then select the positive root because π/5 is in the first quadrant.

\nopagebreak[4]{\small\itshape Answers: see “Exercise 4.3: Elementary functions” under Answers at the end.}
\end{keepsection}
\end{keepsection}

\begin{keepsection}
\Needspace{8\baselineskip}
\section{Self-Study 4.4: Maps}

\Needspace{5\baselineskip}
Vertex form reveals the range and the restriction needed for a quadratic to become one-to-one.

\textbf{Connected topic:} Injectivity and inverses (§3.2) and elementary functions (§3.5); Composition and domains (§3.4).

\begin{keepsection}
\Needspace{7\baselineskip}
\subsection{4.4a: Classify f, g, and h; modify f and g}

\Needspace{4\baselineskip}
Try first: Complete the square for f, solve g for its input, and inspect the coordinate relation defining h.

\nopagebreak[4]{\small\itshape Answers: see “Self-Study 4.4: Maps” under Answers at the end.}
\end{keepsection}

\begin{keepsection}
\Needspace{7\baselineskip}
\subsection{4.4b(i): Form f\ensuremath{\circ}e}

\Needspace{4\baselineskip}
Try first: Since e(x)=e\textsuperscript{x} is always positive, substitute it into f and use the range of the downward parabola.

\nopagebreak[4]{\small\itshape Answers: see “Self-Study 4.4: Maps” under Answers at the end.}
\end{keepsection}

\begin{keepsection}
\Needspace{6\baselineskip}
\subsection{4.4b(ii): Form g\ensuremath{\circ}f}

\Needspace{3\baselineskip}
Try first: The inner output must satisfy f(x)\ensuremath{\geq}0 because g accepts only non-negative inputs.

\nopagebreak[4]{\small\itshape Answers: see “Self-Study 4.4: Maps” under Answers at the end.}
\end{keepsection}

\begin{keepsection}
\Needspace{6\baselineskip}
\subsection{4.4b(iii): Form g\ensuremath{\circ}e}

\Needspace{3\baselineskip}
Try first: Check the exponential output against g’s non-negative domain.

\nopagebreak[4]{\small\itshape Answers: see “Self-Study 4.4: Maps” under Answers at the end.}
\end{keepsection}

\begin{keepsection}
\Needspace{7\baselineskip}
\subsection{4.4b(iv): Form g\ensuremath{\circ}g}

\Needspace{4\baselineskip}
Try first: The first g maps [0,\ensuremath{\infty}) into (0,1], which is contained in the domain [0,\ensuremath{\infty}) of the second g.

\nopagebreak[4]{\small\itshape Answers: see “Self-Study 4.4: Maps” under Answers at the end.}
\end{keepsection}
\end{keepsection}

\begin{keepsection}
\Needspace{8\baselineskip}
\section{Self-Study 4.5: Maps}

\Needspace{5\baselineskip}
Use the definitions directly: equality of composite values, existence of preimages, and the order of quantifiers.

\textbf{Connected topic:} Injectivity and surjectivity (§3.2).

\begin{keepsection}
\Needspace{6\baselineskip}
\subsection{4.5i: If g\ensuremath{\circ}f is injective, show f is injective}

\Needspace{3\baselineskip}
Try first: Assume f sends two inputs to the same value, then apply g.

\nopagebreak[4]{\small\itshape Answers: see “Self-Study 4.5: Maps” under Answers at the end.}
\end{keepsection}

\begin{keepsection}
\Needspace{7\baselineskip}
\subsection{4.5ii: If g\ensuremath{\circ}f is surjective, show g is surjective}

\Needspace{4\baselineskip}
Try first: Take an arbitrary element of Z and use surjectivity of the composite to produce a preimage.

\nopagebreak[4]{\small\itshape Answers: see “Self-Study 4.5: Maps” under Answers at the end.}
\end{keepsection}

\begin{keepsection}
\Needspace{6\baselineskip}
\subsection{4.5iii: If f and g are surjective, show g\ensuremath{\circ}f is surjective}

\Needspace{3\baselineskip}
Try first: For each target z, first choose a g-preimage y, then choose an f-preimage x of y.

\nopagebreak[4]{\small\itshape Answers: see “Self-Study 4.5: Maps” under Answers at the end.}
\end{keepsection}

\begin{keepsection}
\Needspace{6\baselineskip}
\subsection{4.5iv: Give an injective composite with non-injective g}

\Needspace{3\baselineskip}
Try first: Make the domain of f a singleton, while g identifies two different middle elements.

\nopagebreak[4]{\small\itshape Answers: see “Self-Study 4.5: Maps” under Answers at the end.}
\end{keepsection}
\end{keepsection}

\begin{keepsection}
\Needspace{8\baselineskip}
\section{Self-Study 4.6: Elementary functions}

\Needspace{5\baselineskip}
Factor the quartic as a quadratic in x\textsuperscript{2}. For the parameterized quadratic, use the discriminant and include degenerate cases.

\textbf{Connected topic:} Polynomials and elementary functions (§3.5); Elementary functions (§3.5).

\begin{keepsection}
\Needspace{6\baselineskip}
\subsection{4.6a: Find and factor the roots}

\Needspace{3\baselineskip}
Try first: Set u=x\textsuperscript{2} for the quartic, then split the parameter cases for g\textsubscript{a}.

\nopagebreak[4]{\small\itshape Answers: see “Self-Study 4.6: Elementary functions” under Answers at the end.}
\end{keepsection}

\begin{keepsection}
\Needspace{7\baselineskip}
\subsection{4.6b: Show cos(3x) is not generally 3cos x}

\Needspace{4\baselineskip}
Try first: A single convenient counterexample disproves a universal equality; then derive the general expression.

\nopagebreak[4]{\small\itshape Answers: see “Self-Study 4.6: Elementary functions” under Answers at the end.}
\end{keepsection}
\end{keepsection}

\begin{keepsection}
\Needspace{8\baselineskip}
\section{Self-Study 4.7: Strictly monotone functions}

\Needspace{5\baselineskip}
Strict monotonicity forbids equal outputs at distinct inputs; injectivity alone does not require order preservation.

\textbf{Connected topic:} Injectivity and inverse functions (§3.2).

\begin{keepsection}
\Needspace{6\baselineskip}
\subsection{4.7a: Increasing implies injective}

\Needspace{3\baselineskip}
Try first: For two distinct inputs, order them and apply the strict inequality.

\nopagebreak[4]{\small\itshape Answers: see “Self-Study 4.7: Strictly monotone functions” under Answers at the end.}
\end{keepsection}

\begin{keepsection}
\Needspace{6\baselineskip}
\subsection{4.7b: Decreasing implies injective}

\Needspace{3\baselineskip}
Try first: Repeat the order argument with inequalities reversed.

\nopagebreak[4]{\small\itshape Answers: see “Self-Study 4.7: Strictly monotone functions” under Answers at the end.}
\end{keepsection}

\begin{keepsection}
\Needspace{6\baselineskip}
\subsection{4.7c: Is every injection strictly monotone?}

\Needspace{3\baselineskip}
Try first: Find an injective rearrangement of the real line that reverses order for one pair.

\nopagebreak[4]{\small\itshape Answers: see “Self-Study 4.7: Strictly monotone functions” under Answers at the end.}
\end{keepsection}

\begin{keepsection}
\Needspace{6\baselineskip}
\subsection{4.7d: Conclude that exp has an inverse}

\Needspace{3\baselineskip}
Try first: Combine strict increase with the known range of the exponential.

\nopagebreak[4]{\small\itshape Answers: see “Self-Study 4.7: Strictly monotone functions” under Answers at the end.}
\end{keepsection}
\end{keepsection}

\begin{keepsection}
\Needspace{7\baselineskip}
\section{Self-Study 4.8: Tangent addition theorems}

\Needspace{4\baselineskip}
Divide the sine and cosine addition identities only when the stated cosines are nonzero.

\textbf{Connected topic:} Trigonometric functions (§3.5).

\begin{keepsection}
\Needspace{7\baselineskip}
\subsection{4.8a: Prove the tangent addition and subtraction formulas}

\Needspace{4\baselineskip}
Try first: Write tan(x\ensuremath{\pm}y)=sin(x\ensuremath{\pm}y)/cos(x\ensuremath{\pm}y), then divide numerator and denominator by cos x cos y.

\nopagebreak[4]{\small\itshape Answers: see “Self-Study 4.8: Tangent addition theorems” under Answers at the end.}
\end{keepsection}

\begin{keepsection}
\Needspace{7\baselineskip}
\subsection{4.8b: Derive the tangent double-angle formula}

\Needspace{4\baselineskip}
Try first: Set y=x in the addition formula, observing the stated conditions keep the tangent defined.

\nopagebreak[4]{\small\itshape Answers: see “Self-Study 4.8: Tangent addition theorems” under Answers at the end.}
\end{keepsection}

\begin{keepsection}
\Needspace{6\baselineskip}
\subsection{4.8c: Compute tan(π/12) exactly}

\Needspace{3\baselineskip}
Try first: Apply the subtraction formula to π/3\ensuremath{-}π/4 and simplify by rationalizing.

\nopagebreak[4]{\small\itshape Answers: see “Self-Study 4.8: Tangent addition theorems” under Answers at the end.}
\end{keepsection}
\end{keepsection}

\begin{keepsection}
\Needspace{10\baselineskip}
\section{Challenge 4.A: Surjective maps}

\Needspace{7\baselineskip}
Let s(k) count maps onto a fixed k-element set. Partition all maps according to their image, then invert the resulting triangular relation.

\textbf{Connected topic:} Map types (§3.2) and counting and binomial coefficients (§2.4); Finite sums and binomial coefficients (§2.4); Binomial coefficients (§2.4); Binomial coefficients and finite sums (§2.4); Induction (§2.5) and map types (§3.2).

\begin{keepsection}
\Needspace{7\baselineskip}
\subsection{4.A setup: Interpret s(k) and the target}

\Needspace{4\baselineskip}
Try first: Relate surjectivity to the size of the image, and identify s(n) with the desired count.

\nopagebreak[4]{\small\itshape Answers: see “Challenge 4.A: Surjective maps” under Answers at the end.}
\end{keepsection}

\begin{keepsection}
\Needspace{6\baselineskip}
\subsection{4.A Step 1: Derive the recurrence for s(k)}

\Needspace{3\baselineskip}
Try first: Partition the k\textsuperscript{m} maps by their exact image size j and solve for the j=k term.

\nopagebreak[4]{\small\itshape Answers: see “Challenge 4.A: Surjective maps” under Answers at the end.}
\end{keepsection}

\begin{keepsection}
\Needspace{7\baselineskip}
\subsection{4.A Step 2: Check the proposed formula for k=1, 2, and 3}

\Needspace{4\baselineskip}
Try first: Use the recurrence to calculate the first values and compare them with the binomial sum.

\nopagebreak[4]{\small\itshape Answers: see “Challenge 4.A: Surjective maps” under Answers at the end.}
\end{keepsection}

\begin{keepsection}
\Needspace{6\baselineskip}
\subsection{4.A Step 3a: Reverse the order of a triangular sum}

\Needspace{3\baselineskip}
Try first: List the index region rather than manipulating terms.

\nopagebreak[4]{\small\itshape Answers: see “Challenge 4.A: Surjective maps” under Answers at the end.}
\end{keepsection}

\begin{keepsection}
\Needspace{6\baselineskip}
\subsection{4.A Step 3b: Prove the binomial coefficient identity}

\Needspace{3\baselineskip}
Try first: Write both sides using factorial definitions.

\nopagebreak[4]{\small\itshape Answers: see “Challenge 4.A: Surjective maps” under Answers at the end.}
\end{keepsection}

\begin{keepsection}
\Needspace{7\baselineskip}
\subsection{4.A Step 3c: Establish the convolution identity}

\Needspace{4\baselineskip}
Try first: After changing order, the coefficient of each i contains a complete alternating binomial sum.

\nopagebreak[4]{\small\itshape Answers: see “Challenge 4.A: Surjective maps” under Answers at the end.}
\end{keepsection}

\begin{keepsection}
\Needspace{7\baselineskip}
\subsection{4.A Step 4: Prove s(k)=g(k) by induction}

\Needspace{4\baselineskip}
Try first: Use k=1 as the base, then subtract the same lower-index terms from the recurrence and convolution identity.

\nopagebreak[4]{\small\itshape Answers: see “Challenge 4.A: Surjective maps” under Answers at the end.}
\end{keepsection}
\end{keepsection}


\clearpage
\section{Answers}
\begin{keepsection}\subsection{Exercise 4.1: Mathematical induction}

\begin{workedbox}{4.1a \ensuremath{\cdot} Prove the finite geometric sum}
\Needspace{11\baselineskip}
\begin{enumerate}
\item Let S\textsubscript{n}=\ensuremath{\Sigma}\textsubscript{k=0}\textsuperscript{n}x\textsuperscript{k}. The claim is S\textsubscript{n}=(1\ensuremath{-}x\textsuperscript{n+1})/(1\ensuremath{-}x), for x\ensuremath{\neq}1.
\item At n=0, S\textsubscript{0}=1 and (1\ensuremath{-}x)/(1\ensuremath{-}x)=1.
\item Assume the formula for n. Add the next term: S\textsubscript{n+1}=S\textsubscript{n}+x\textsuperscript{n+1}=(1\ensuremath{-}x\textsuperscript{n+1})/(1\ensuremath{-}x)+x\textsuperscript{n+1}.
\item Put the terms over the same denominator: [1\ensuremath{-}x\textsuperscript{n+1}+x\textsuperscript{n+1}(1\ensuremath{-}x)]/(1\ensuremath{-}x)=(1\ensuremath{-}x\textsuperscript{n+2})/(1\ensuremath{-}x).
\item That is the claimed formula with n+1. The denominator is nonzero because x\ensuremath{\neq}1; the base and step prove the result for all n\ensuremath{\geq}0.
\item \begin{equation*}\sum_{\mathrm{k}=0}^{\mathrm{n}}\, \mathrm{x}^{\mathrm{k}}\, =\, (1-\mathrm{x}^{\mathrm{n}+1})/(1-\mathrm{x})\end{equation*}
\end{enumerate}
\end{workedbox}

\end{keepsection}
\begin{keepsection}\subsection{Exercise 4.1: Mathematical induction}

\begin{workedbox}{4.1b \ensuremath{\cdot} Prove the reciprocal-square-root bound}
\Needspace{8\baselineskip}
\begin{enumerate}
\item For 1\ensuremath{\leq}k\ensuremath{\leq}n, \ensuremath{\sqrt{\mathrm{k}}}\ensuremath{\leq}\ensuremath{\sqrt{\mathrm{n}}}, so 1/\ensuremath{\sqrt{\mathrm{k}}}\ensuremath{\geq}1/\ensuremath{\sqrt{\mathrm{n}}}.
\item Summing these n inequalities gives \ensuremath{\Sigma}\textsubscript{k=1}\textsuperscript{n}1/\ensuremath{\sqrt{\mathrm{k}}}\ensuremath{\geq}n/\ensuremath{\sqrt{\mathrm{n}}}=\ensuremath{\sqrt{\mathrm{n}}}.
\item For n\ensuremath{\geq}2 at least one inequality is strict (for example k=1k=1\textsuperscript{n}1/\ensuremath{\sqrt{\mathrm{k}}}\ensuremath{>}\ensuremath{\sqrt{\mathrm{n}}}.
\item This proves the stated claim directly, and also explains why the inequality is strict.
\item \begin{equation*}\sum_{\mathrm{k}=1}^{\mathrm{n}}\, 1/\sqrt{\mathrm{k}}\, >\, \sqrt{\mathrm{n}}\end{equation*}
\end{enumerate}
\end{workedbox}

\end{keepsection}
\begin{keepsection}\subsection{Exercise 4.2: Maps}

\begin{workedbox}{4.2a \ensuremath{\cdot} Is a composition of injections injective?}
\Needspace{8\baselineskip}
\begin{enumerate}
\item Suppose (g\ensuremath{\circ}f)(x\textsubscript{1})=(g\ensuremath{\circ}f)(x\textsubscript{2}). This says g(f(x\textsubscript{1}))=g(f(x\textsubscript{2})).
\item Since g is injective, f(x\textsubscript{1})=f(x\textsubscript{2}). Since f is injective, x\textsubscript{1}=x\textsubscript{2}.
\item Thus g\ensuremath{\circ}f is injective. The codomains need not be equal: the key is that each map is injective on the outputs that occur.
\item \begin{equation*}(\mathrm{g}\circ \mathrm{f})(\mathrm{x}_{1})=(\mathrm{g}\circ \mathrm{f})(\mathrm{x}_{2})\, \Rightarrow \, \mathrm{x}_{1}=\mathrm{x}_{2}\end{equation*}
\end{enumerate}
\end{workedbox}

\end{keepsection}
\begin{keepsection}\subsection{Exercise 4.2: Maps}

\begin{workedbox}{4.2b(i) \ensuremath{\cdot} Count all maps M\ensuremath{\to}N}
\Needspace{9\baselineskip}
\begin{enumerate}
\item For each input x\ensuremath{\in}M, a map may assign any of the n elements of N.
\item There are n choices for each of m independent inputs, so the product rule gives n\ensuremath{\cdot}n\ensuremath{\cdot}\ensuremath{\cdot}\ensuremath{\cdot}n=n\textsuperscript{m} maps.
\item If n=0 and m\ensuremath{>}0, this gives 0 maps, as expected; for m=0 the unique empty map is counted by 0\textsuperscript{0}=1 under the combinatorial convention.
\item \begin{equation*}\#(\mathrm{M}\to \mathrm{N})=\mathrm{n}^{\mathrm{m}}\end{equation*}
\end{enumerate}
\end{workedbox}

\end{keepsection}
\begin{keepsection}\subsection{Exercise 4.2: Maps}

\begin{workedbox}{4.2b(ii) \ensuremath{\cdot} Count bijections M\ensuremath{\to}N}
\Needspace{8\baselineskip}
\begin{enumerate}
\item If m\ensuremath{\neq}n, no bijection exists, so the count is 0.
\item If m=n, list the m inputs in any fixed order. The first image has m choices, the next m\ensuremath{-}1, and so on to 1.
\item The product is m!, including 0!=1 for the empty sets.
\item \begin{equation*}\#\{\text{bijective}\, \text{maps}\, \mathrm{M}\to \mathrm{N}\}=\mathrm{m}!\, \text{when}\, \mathrm{m}=\mathrm{n};\, \text{otherwise}\, 0\end{equation*}
\end{enumerate}
\end{workedbox}

\end{keepsection}
\begin{keepsection}\subsection{Exercise 4.2: Maps}

\begin{workedbox}{4.2b(iii) \ensuremath{\cdot} Count injective maps M\ensuremath{\to}N}
\Needspace{8\baselineskip}
\begin{enumerate}
\item If m\ensuremath{>}n, the pigeonhole principle makes an injection impossible; the count is 0.
\item For m\ensuremath{\leq}n, order the inputs. Choose distinct images in sequence: n choices, then n\ensuremath{-}1, down to n\ensuremath{-}m+1.
\item Therefore the count is n(n\ensuremath{-}1)…(n\ensuremath{-}m+1)=n!/(n\ensuremath{-}m)!. For m=0 this is 1.
\item \begin{equation*}\#\{\text{injective}\, \text{maps}\, \mathrm{M}\to \mathrm{N}\}=\mathrm{n}!/(\mathrm{n}-\mathrm{m})!\, \text{when}\, \mathrm{m}\leq \mathrm{n};\, \text{otherwise}\, 0\end{equation*}
\end{enumerate}
\end{workedbox}

\end{keepsection}
\begin{keepsection}\subsection{Exercise 4.2: Maps}

\begin{workedbox}{4.2b(iv) \ensuremath{\cdot} Count surjective maps M\ensuremath{\to}N}
\Needspace{10\baselineskip}
\begin{enumerate}
\item Start with all n\textsuperscript{m} maps. For a chosen set of j target elements forbidden from the image, the remaining n\ensuremath{-}j targets give (n\ensuremath{-}j)\textsuperscript{m} maps.
\item There are C(n,j) ways to choose those j omitted targets. Inclusion–exclusion alternates the correction by the number omitted.
\item The number of onto maps is \ensuremath{\Sigma}\textsubscript{j=0}\textsuperscript{n}(\ensuremath{-}1)\textsuperscript{j} C(n,j)(n\ensuremath{-}j)\textsuperscript{m}. This includes the empty-set cases when interpreted with the usual combinatorial convention.
\item \begin{equation*}\#\{\text{onto}\, \text{maps}\, \mathrm{M}\to \mathrm{N}\}=\sum_{\mathrm{j}=0}^{\mathrm{n}}(-1)^{\mathrm{j}}\, \mathrm{C}(\mathrm{n},\mathrm{j})(\mathrm{n}-\mathrm{j})^{\mathrm{m}}\end{equation*}
\end{enumerate}
\end{workedbox}

\end{keepsection}
\begin{keepsection}\subsection{Exercise 4.3: Elementary functions}

\begin{workedbox}{4.3a \ensuremath{\cdot} Prove the circle equation}
\Needspace{7\baselineskip}
\begin{enumerate}
\item The cosine difference formula is cos(a\ensuremath{-}b)=cos a cos b+sin a sin b.
\item Set a=b=x. The left side is cos 0=1, and the right side is cos\textsuperscript{2}x+sin\textsuperscript{2}x.
\item Hence cos\textsuperscript{2}x+sin\textsuperscript{2}x=1 for every real x.
\item \begin{equation*}\cos ^{2}\mathrm{x}+\sin ^{2}\mathrm{x}=1\end{equation*}
\end{enumerate}
\end{workedbox}

\end{keepsection}
\begin{keepsection}\subsection{Exercise 4.3: Elementary functions}

\begin{workedbox}{4.3b \ensuremath{\cdot} Prove the cosine double-angle identity}
\Needspace{5\baselineskip}
\begin{enumerate}
\item The addition formula gives cos(x+x)=cos\textsuperscript{2}x\ensuremath{-}sin\textsuperscript{2}x.
\item Using cos\textsuperscript{2}x=1\ensuremath{-}sin\textsuperscript{2}x from the circle equation, cos(2x)=1\ensuremath{-}sin\textsuperscript{2}x\ensuremath{-}sin\textsuperscript{2}x=1\ensuremath{-}2sin\textsuperscript{2}x.
\item \begin{equation*}\cos (2\mathrm{x})=1-2\sin ^{2}\mathrm{x}\end{equation*}
\end{enumerate}
\end{workedbox}

\end{keepsection}
\begin{keepsection}\subsection{Exercise 4.3: Elementary functions}

\begin{workedbox}{4.3c \ensuremath{\cdot} Evaluate cos(arcsin x)}
\Needspace{8\baselineskip}
\begin{enumerate}
\item By definition, θ=arcsin x means sin θ=x, with θ\ensuremath{\in}[\ensuremath{-}π/2,π/2].
\item The circle identity gives cos\textsuperscript{2}θ=1\ensuremath{-}sin\textsuperscript{2}θ=1\ensuremath{-}x\textsuperscript{2}. On this interval cos θ\ensuremath{\geq}0, so take the non-negative square root.
\item Therefore cos(arcsin x)=\ensuremath{\sqrt{(1-\mathrm{x}²)}}, valid for \ensuremath{-}1\ensuremath{\leq}x\ensuremath{\leq}1.
\item \begin{equation*}\cos (\arcsin \, \mathrm{x})=\sqrt{(1-\mathrm{x}²)},\, -1\leq \mathrm{x}\leq 1\end{equation*}
\end{enumerate}
\end{workedbox}

\end{keepsection}
\begin{keepsection}\subsection{Exercise 4.3: Elementary functions}

\begin{workedbox}{4.3d \ensuremath{\cdot} Find cos(π/5) exactly}
\Needspace{7\baselineskip}
\begin{enumerate}
\item The given value is sin(π/10)=(\ensuremath{\sqrt{5}}\ensuremath{-}1)/4. Its square is (3\ensuremath{-}\ensuremath{\sqrt{5}})/8.
\item Apply cos(2x)=1\ensuremath{-}2sin\textsuperscript{2}x with x=π/10: cos(π/5)=1\ensuremath{-}2(3\ensuremath{-}\ensuremath{\sqrt{5}})/8=(1+\ensuremath{\sqrt{5}})/4.
\item The result is positive, as expected for π/5\ensuremath{\in}(0,π/2).
\item \begin{equation*}\cos (π/5)=(1+\sqrt{5})/4\end{equation*}
\end{enumerate}
\end{workedbox}

\end{keepsection}
\begin{keepsection}\subsection{Self-Study 4.4: Maps}

\begin{workedbox}{4.4a \ensuremath{\cdot} Classify f, g, and h; modify f and g}
\Needspace{15\baselineskip}
\begin{enumerate}
\item For f(x)=\ensuremath{-}x\textsuperscript{2}+6x\ensuremath{-}5, complete the square: f(x)=\ensuremath{-}(x\ensuremath{-}3)\textsuperscript{2}+4. Its range is (\ensuremath{-}\ensuremath{\infty},4]; it is not injective on \ensuremath{\mathbb{R}} because points equally spaced about 3 have equal values.
\item Restrict f to [3,\ensuremath{\infty}) and use codomain (\ensuremath{-}\ensuremath{\infty},4]. It is then strictly decreasing and onto that codomain. Solving y=4\ensuremath{-}(x\ensuremath{-}3)\textsuperscript{2} on this branch gives f\textsuperscript{\ensuremath{-}1}(y)=3+\ensuremath{\sqrt{(4-\mathrm{y})}}, y\ensuremath{\leq}4.
\item For g(x)=1/(1+x), x\ensuremath{\geq}0, the values are (0,1]. It is strictly decreasing and injective, but not onto \ensuremath{\mathbb{R}}. With codomain (0,1], it is bijective; solving y=1/(1+x) gives g\textsuperscript{\ensuremath{-}1}(y)=1/y\ensuremath{-}1.
\item For h(x,y)=(x,x+y,y), equality of outputs gives equality of the first and third coordinates, hence x and y agree: h is injective. Its image is the plane \{(a,b,c):b=a+c\}, so it is not onto \ensuremath{\mathbb{R}}\textsuperscript{3} and is not bijective.
\item \begin{equation*}\mathrm{f}(\mathrm{x})=-(\mathrm{x}-3)^{2}+4;\, \mathrm{f}^{-1}(\mathrm{y})=3+\sqrt{(4-\mathrm{y})}\end{equation*}
\item \begin{equation*}\mathrm{g}^{-1}(\mathrm{y})=1/\mathrm{y}-1\end{equation*}
\end{enumerate}
\end{workedbox}

\end{keepsection}
\begin{keepsection}\subsection{Self-Study 4.4: Maps}

\begin{workedbox}{4.4b(i) \ensuremath{\cdot} Form f\ensuremath{\circ}e}
\Needspace{8\baselineskip}
\begin{enumerate}
\item Every real x is allowed: e\textsuperscript{x}\ensuremath{\in}\ensuremath{\mathbb{R}}, the domain of f.
\item (f\ensuremath{\circ}e)(x)=\ensuremath{-}(e\textsuperscript{x})\textsuperscript{2}+6e\textsuperscript{x}\ensuremath{-}5=\ensuremath{-}e\textsuperscript{2x}+6e\textsuperscript{x}\ensuremath{-}5.
\item Writing t=e\textsuperscript{x}\ensuremath{>}0 gives 4\ensuremath{-}(t\ensuremath{-}3)\textsuperscript{2}. Its maximum 4 occurs when t=3; as t\ensuremath{\to}0\textsuperscript{+} or t\ensuremath{\to}\ensuremath{\infty} the value tends without bound downward. The range is (\ensuremath{-}\ensuremath{\infty},4].
\item \begin{equation*}(\mathrm{f}\circ \mathrm{e})(\mathrm{x})=-\mathrm{e}^{2\mathrm{x}}+6\mathrm{e}^{\mathrm{x}}-5;\, \text{range}\, (-\infty ,4]\end{equation*}
\end{enumerate}
\end{workedbox}

\end{keepsection}
\begin{keepsection}\subsection{Self-Study 4.4: Maps}

\begin{workedbox}{4.4b(ii) \ensuremath{\cdot} Form g\ensuremath{\circ}f}
\Needspace{7\baselineskip}
\begin{enumerate}
\item Require 4\ensuremath{-}(x\ensuremath{-}3)\textsuperscript{2}\ensuremath{\geq}0, so (x\ensuremath{-}3)\textsuperscript{2}\ensuremath{\leq}4 and 1\ensuremath{\leq}x\ensuremath{\leq}5. This is the composite domain.
\item On it, f(x) ranges over [0,4], and (g\ensuremath{\circ}f)(x)=1/(1+f(x))=1/[5\ensuremath{-}(x\ensuremath{-}3)\textsuperscript{2}].
\item The denominator runs from 1 to 5; reciprocation reverses the order. The range is [1/5,1].
\item \begin{equation*}(\mathrm{g}\circ \mathrm{f})(\mathrm{x})=1/[5-(\mathrm{x}-3)^{2}],\, \text{domain}\, [1,5],\, \text{range}\, [1/5,1]\end{equation*}
\end{enumerate}
\end{workedbox}

\end{keepsection}
\begin{keepsection}\subsection{Self-Study 4.4: Maps}

\begin{workedbox}{4.4b(iii) \ensuremath{\cdot} Form g\ensuremath{\circ}e}
\Needspace{8\baselineskip}
\begin{enumerate}
\item For all x\ensuremath{\in}\ensuremath{\mathbb{R}}, e\textsuperscript{x}\ensuremath{>}0, so the composition is defined everywhere.
\item (g\ensuremath{\circ}e)(x)=1/(1+e\textsuperscript{x}).
\item As e\textsuperscript{x} ranges over (0,\ensuremath{\infty}), this expression decreases through (0,1). It approaches 1 and 0 but reaches neither, so the range is (0,1).
\item \begin{equation*}(\mathrm{g}\circ \mathrm{e})(\mathrm{x})=1/(1+\mathrm{e}^{\mathrm{x}}),\, \text{range}\, (0,1)\end{equation*}
\end{enumerate}
\end{workedbox}

\end{keepsection}
\begin{keepsection}\subsection{Self-Study 4.4: Maps}

\begin{workedbox}{4.4b(iv) \ensuremath{\cdot} Form g\ensuremath{\circ}g}
\Needspace{7\baselineskip}
\begin{enumerate}
\item The composite domain is [0,\ensuremath{\infty}), since g(x)\ensuremath{\geq}0 for every allowed x.
\item (g\ensuremath{\circ}g)(x)=1/[1+1/(1+x)]=(1+x)/(2+x).
\item This increases from 1/2 at x=0 and approaches 1 without reaching it. The range is [1/2,1).
\item \begin{equation*}(\mathrm{g}\circ \mathrm{g})(\mathrm{x})=(1+\mathrm{x})/(2+\mathrm{x}),\, \text{range}\, [1/2,1)\end{equation*}
\end{enumerate}
\end{workedbox}

\end{keepsection}
\begin{keepsection}\subsection{Self-Study 4.5: Maps}

\begin{workedbox}{4.5i \ensuremath{\cdot} If g\ensuremath{\circ}f is injective, show f is injective}
\Needspace{5\baselineskip}
\begin{enumerate}
\item Suppose f(x\textsubscript{1})=f(x\textsubscript{2}). Applying g gives (g\ensuremath{\circ}f)(x\textsubscript{1})=(g\ensuremath{\circ}f)(x\textsubscript{2}).
\item Injectivity of g\ensuremath{\circ}f implies x\textsubscript{1}=x\textsubscript{2}. Thus f is injective.
\item \begin{equation*}\mathrm{g}\circ \mathrm{f}\, \text{injective}\, \Rightarrow \, \mathrm{f}\, \text{injective}\end{equation*}
\end{enumerate}
\end{workedbox}

\end{keepsection}
\begin{keepsection}\subsection{Self-Study 4.5: Maps}

\begin{workedbox}{4.5ii \ensuremath{\cdot} If g\ensuremath{\circ}f is surjective, show g is surjective}
\Needspace{6\baselineskip}
\begin{enumerate}
\item Let z\ensuremath{\in}Z. Since g\ensuremath{\circ}f is onto Z, some x\ensuremath{\in}X satisfies g(f(x))=z.
\item The element y=f(x) belongs to Y and g(y)=z. Every z has a preimage under g, so g is surjective.
\item \begin{equation*}\mathrm{g}\circ \mathrm{f}\, \text{surjective}\, \Rightarrow \, \mathrm{g}\, \text{surjective}\end{equation*}
\end{enumerate}
\end{workedbox}

\end{keepsection}
\begin{keepsection}\subsection{Self-Study 4.5: Maps}

\begin{workedbox}{4.5iii \ensuremath{\cdot} If f and g are surjective, show g\ensuremath{\circ}f is surjective}
\Needspace{7\baselineskip}
\begin{enumerate}
\item Let z\ensuremath{\in}Z. Surjectivity of g gives y\ensuremath{\in}Y with g(y)=z.
\item Surjectivity of f gives x\ensuremath{\in}X with f(x)=y. Therefore g(f(x))=z.
\item Every z is reached by g\ensuremath{\circ}f, proving surjectivity.
\item \begin{equation*}\mathrm{f}\, \text{and}\, \mathrm{g}\, \text{surjective}\, \Rightarrow \, \mathrm{g}\circ \mathrm{f}\, \text{surjective}\end{equation*}
\end{enumerate}
\end{workedbox}

\end{keepsection}
\begin{keepsection}\subsection{Self-Study 4.5: Maps}

\begin{workedbox}{4.5iv \ensuremath{\cdot} Give an injective composite with non-injective g}
\Needspace{8\baselineskip}
\begin{enumerate}
\item Let X=\{a\}, Y=\{0,1\}, Z=\{z\}. Define f(a)=0 and define g(0)=g(1)=z.
\item The map g is not injective because 0\ensuremath{\neq}1 but g(0)=g(1).
\item The composite X\ensuremath{\to}Z sends the only input a to z. A map on a singleton is injective, so g\ensuremath{\circ}f is injective.
\item \begin{equation*}\mathrm{X}=\{\mathrm{a}\},\, \mathrm{Y}=\{0,1\},\, \mathrm{Z}=\{\mathrm{z}\};\, \mathrm{f}(\mathrm{a})=0,\, \mathrm{g}(0)=\mathrm{g}(1)=\mathrm{z}\end{equation*}
\end{enumerate}
\end{workedbox}

\end{keepsection}
\begin{keepsection}\subsection{Self-Study 4.6: Elementary functions}

\begin{workedbox}{4.6a \ensuremath{\cdot} Find and factor the roots}
\Needspace{14\baselineskip}
\begin{enumerate}
\item For f, divide by 5 and set u=x\textsuperscript{2}: u\textsuperscript{2}+4u\ensuremath{-}5=(u\ensuremath{-}1)(u+5). Thus f(x)=5(x\textsuperscript{2}\ensuremath{-}1)(x\textsuperscript{2}+5).
\item The real roots are x=\ensuremath{\pm}1. Over \ensuremath{\mathbb{C}} the complete linear factorization is 5(x\ensuremath{-}1)(x+1)(x\ensuremath{-}i\ensuremath{\sqrt{5}})(x+i\ensuremath{\sqrt{5}}).
\item For g\textsubscript{a}(x)=ax\textsuperscript{2}+2x+1: if a=0 it is 2x+1, with one root \ensuremath{-}1/2. If a\ensuremath{\neq}0, its discriminant is Δ=4\ensuremath{-}4a.
\item When a\ensuremath{<}1 and a\ensuremath{\neq}0 there are two distinct real roots r\ensuremath{\pm}=(\ensuremath{-}1\ensuremath{\pm}\ensuremath{\sqrt{(1-\mathrm{a})}})/a, so g\textsubscript{a}=a(x\ensuremath{-}r\textsubscript{+})(x\ensuremath{-}r\textsubscript{\ensuremath{-}}). At a=1 it is (x+1)\textsuperscript{2}.
\item When a\ensuremath{>}1 there are no real roots; over \ensuremath{\mathbb{C}} its factors are a(x\ensuremath{-}r\textsubscript{+})(x\ensuremath{-}r\textsubscript{\ensuremath{-}}), where r\ensuremath{\pm}=(\ensuremath{-}1\ensuremath{\pm}i\ensuremath{\sqrt{(\mathrm{a}-1)}})/a.
\item \begin{equation*}\mathrm{f}(\mathrm{x})=5(\mathrm{x}-1)(\mathrm{x}+1)(\mathrm{x}-\mathrm{i}\sqrt{5})(\mathrm{x}+\mathrm{i}\sqrt{5})\end{equation*}
\item \begin{equation*}Δ(\mathrm{g}_{\mathrm{a}})=4(1-\mathrm{a})\end{equation*}
\end{enumerate}
\end{workedbox}

\end{keepsection}
\begin{keepsection}\subsection{Self-Study 4.6: Elementary functions}

\begin{workedbox}{4.6b \ensuremath{\cdot} Show cos(3x) is not generally 3cos x}
\Needspace{7\baselineskip}
\begin{enumerate}
\item Take x=0: cos(3\ensuremath{\cdot}0)=1 while 3cos 0=3, so they are not equal in general.
\item Use cos(2x+x)=cos(2x)cos x\ensuremath{-}sin(2x)sin x. Substitute cos2x=cos\textsuperscript{2}x\ensuremath{-}sin\textsuperscript{2}x and sin2x=2sin x cos x.
\item This gives cos3x=cos\textsuperscript{3}x\ensuremath{-}3sin\textsuperscript{2}x cos x=cos\textsuperscript{3}x\ensuremath{-}3(1\ensuremath{-}cos\textsuperscript{2}x)cos x=4cos\textsuperscript{3}x\ensuremath{-}3cos x.
\item \begin{equation*}\cos (3\mathrm{x})=4\cos ^{3}\mathrm{x}-3\cos \, \mathrm{x}\end{equation*}
\end{enumerate}
\end{workedbox}

\end{keepsection}
\begin{keepsection}\subsection{Self-Study 4.7: Strictly monotone functions}

\begin{workedbox}{4.7a \ensuremath{\cdot} Increasing implies injective}
\Needspace{5\baselineskip}
\begin{enumerate}
\item Take x\ensuremath{\neq}y. Either x
\item If y
\item \begin{equation*}\mathrm{x}\end{equation*}
\end{enumerate}
\end{workedbox}

\end{keepsection}
\begin{keepsection}\subsection{Self-Study 4.7: Strictly monotone functions}

\begin{workedbox}{4.7b \ensuremath{\cdot} Decreasing implies injective}
\Needspace{5\baselineskip}
\begin{enumerate}
\item Take x\ensuremath{\neq}y. If xf(y); if yf(x).
\item In either case the outputs differ, so f is injective.
\item \begin{equation*}\mathrm{x}\mathrm{f}(\mathrm{y})\, \Rightarrow \, \mathrm{f}(\mathrm{x})≠\mathrm{f}(\mathrm{y})\end{equation*}
\end{enumerate}
\end{workedbox}

\end{keepsection}
\begin{keepsection}\subsection{Self-Study 4.7: Strictly monotone functions}

\begin{workedbox}{4.7c \ensuremath{\cdot} Is every injection strictly monotone?}
\Needspace{9\baselineskip}
\begin{enumerate}
\item Define f:\ensuremath{\mathbb{R}}\ensuremath{\to}\ensuremath{\mathbb{R}} by swapping 0 and 1 and leaving every other real number fixed. This is a bijection, hence injective.
\item Since 0\ensuremath{<}1/2 but f(0)=1\ensuremath{>}f(1/2)=1/2, it is not increasing. Since 1/2\ensuremath{<}1 but f(1/2)=1/2\ensuremath{>}f(1)=0, it is not decreasing.
\item Thus injectivity does not imply strict monotonicity.
\item \begin{equation*}\mathrm{f}(0)=1>\mathrm{f}(1/2)=1/2,\, \mathrm{f}(1/2)=1/2>\mathrm{f}(1)=0\end{equation*}
\end{enumerate}
\end{workedbox}

\end{keepsection}
\begin{keepsection}\subsection{Self-Study 4.7: Strictly monotone functions}

\begin{workedbox}{4.7d \ensuremath{\cdot} Conclude that exp has an inverse}
\Needspace{7\baselineskip}
\begin{enumerate}
\item The exponential is strictly increasing on \ensuremath{\mathbb{R}}, so by part (a) it is injective.
\item Its range is (0,\ensuremath{\infty}), so exp:\ensuremath{\mathbb{R}}\ensuremath{\to}(0,\ensuremath{\infty}) is also surjective.
\item It is therefore bijective and has an inverse, the natural logarithm ln:(0,\ensuremath{\infty})\ensuremath{\to}\ensuremath{\mathbb{R}}.
\item \begin{equation*}\exp :\mathbb{R}\to (0,\infty )\, \text{is}\, \text{bijective};\, \exp ^{-1}=\ln \end{equation*}
\end{enumerate}
\end{workedbox}

\end{keepsection}
\begin{keepsection}\subsection{Self-Study 4.8: Tangent addition theorems}

\begin{workedbox}{4.8a \ensuremath{\cdot} Prove the tangent addition and subtraction formulas}
\Needspace{10\baselineskip}
\begin{enumerate}
\item The addition identities give sin(x\ensuremath{\pm}y)=sin x cos y\ensuremath{\pm}cos x sin y and cos(x\ensuremath{\pm}y)=cos x cos y\ensuremath{\mp}sin x sin y.
\item The hypotheses ensure cos x, cos y, and cos(x\ensuremath{\pm}y) are nonzero, so division is valid. Divide top and bottom by cos x cos y.
\item The result is tan(x\ensuremath{\pm}y)=(tan x\ensuremath{\pm}tan y)/(1\ensuremath{\mp}tan x tan y). For the plus case use matching upper/lower signs; for the minus case use opposite signs.
\item \begin{equation*}\tan (\mathrm{x}\pm \mathrm{y})=(\tan \, \mathrm{x}\pm \tan \, \mathrm{y})/(1\mp \tan \, \mathrm{x}\, \tan \, \mathrm{y})\end{equation*}
\end{enumerate}
\end{workedbox}

\end{keepsection}
\begin{keepsection}\subsection{Self-Study 4.8: Tangent addition theorems}

\begin{workedbox}{4.8b \ensuremath{\cdot} Derive the tangent double-angle formula}
\Needspace{5\baselineskip}
\begin{enumerate}
\item Use tan(x+x)=(tan x+tan x)/(1\ensuremath{-}tan\textsuperscript{2}x).
\item Thus tan(2x)=2tan x/(1\ensuremath{-}tan\textsuperscript{2}x), whenever x and 2x avoid the excluded tangent poles.
\item \begin{equation*}\tan (2\mathrm{x})=2\tan \, \mathrm{x}/(1-\tan ^{2}\mathrm{x})\end{equation*}
\end{enumerate}
\end{workedbox}

\end{keepsection}
\begin{keepsection}\subsection{Self-Study 4.8: Tangent addition theorems}

\begin{workedbox}{4.8c \ensuremath{\cdot} Compute tan(π/12) exactly}
\Needspace{8\baselineskip}
\begin{enumerate}
\item Use tan(π/3)=\ensuremath{\sqrt{3}} and tan(π/4)=1. Then tan(π/12)=(\ensuremath{\sqrt{3}}\ensuremath{-}1)/(1+\ensuremath{\sqrt{3}}).
\item Multiply numerator and denominator by \ensuremath{\sqrt{3}}\ensuremath{-}1: the denominator becomes 3\ensuremath{-}1=2 and the numerator becomes (\ensuremath{\sqrt{3}}\ensuremath{-}1)\textsuperscript{2}=4\ensuremath{-}2\ensuremath{\sqrt{3}}.
\item Therefore tan(π/12)=2\ensuremath{-}\ensuremath{\sqrt{3}}, a positive value as expected for an angle of 15°.
\item \begin{equation*}\tan (π/12)=2-\sqrt{3}\end{equation*}
\end{enumerate}
\end{workedbox}

\end{keepsection}
\begin{keepsection}\subsection{Challenge 4.A: Surjective maps}

\begin{workedbox}{4.A setup \ensuremath{\cdot} Interpret s(k) and the target}
\Needspace{9\baselineskip}
\begin{enumerate}
\item For a fixed k-element set A, s(k) is the number of maps f:M\ensuremath{\to}A whose image is all of A.
\item A map M\ensuremath{\to}N is surjective exactly when its image has n elements, so s(n) counts the onto maps from M to N.
\item Assume m=|M|\ensuremath{\geq}1 for the stated Step 1 recurrence. If m=0, there is one onto map when n=0 and none when n\ensuremath{>}0; that case is separate from this recurrence.
\item \begin{equation*}\mathrm{f}\, \text{is}\, \text{onto}\, \Leftrightarrow \, |\mathrm{f}(\mathrm{M})|=|\mathrm{N}|=\mathrm{n}\end{equation*}
\end{enumerate}
\end{workedbox}

\end{keepsection}
\begin{keepsection}\subsection{Challenge 4.A: Surjective maps}

\begin{workedbox}{4.A Step 1 \ensuremath{\cdot} Derive the recurrence for s(k)}
\Needspace{8\baselineskip}
\begin{enumerate}
\item There are k\textsuperscript{m} total maps M\ensuremath{\to}A. If the image has size j, first choose its j-element image B in C(k,j) ways, then choose a map onto B in s(j) ways.
\item Summing over image sizes gives k\textsuperscript{m}=\ensuremath{\Sigma}\textsubscript{j=1}\textsuperscript{k} C(k,j)s(j). There is no empty image because m\ensuremath{\geq}1.
\item The j=k term is s(k), since C(k,k)=1. Rearranging gives s(k)=k\textsuperscript{m}\ensuremath{-}\ensuremath{\Sigma}\textsubscript{j=1}\textsuperscript{k\ensuremath{-}1} C(k,j)s(j).
\item \begin{equation*}\mathrm{k}^{\mathrm{m}}=\sum_{\mathrm{j}=1}^{\mathrm{k}}\, \mathrm{C}(\mathrm{k},\mathrm{j})\mathrm{s}(\mathrm{j});\, \mathrm{s}(\mathrm{k})=\mathrm{k}^{\mathrm{m}}-\sum_{\mathrm{j}=1}^{\mathrm{k}-1}\, \mathrm{C}(\mathrm{k},\mathrm{j})\mathrm{s}(\mathrm{j})\end{equation*}
\end{enumerate}
\end{workedbox}

\end{keepsection}
\begin{keepsection}\subsection{Challenge 4.A: Surjective maps}

\begin{workedbox}{4.A Step 2 \ensuremath{\cdot} Check the proposed formula for k=1, 2, and 3}
\Needspace{8\baselineskip}
\begin{enumerate}
\item For k=1, s(1)=1\textsuperscript{m}=1. The proposed sum is C(1,1)1\textsuperscript{m}=1.
\item For k=2, s(2)=2\textsuperscript{m}\ensuremath{-}2s(1)=2\textsuperscript{m}\ensuremath{-}2. The guess gives C(2,1)(\ensuremath{-}1)\textsuperscript{1}1\textsuperscript{m}+C(2,2)2\textsuperscript{m}=2\textsuperscript{m}\ensuremath{-}2.
\item For k=3, s(3)=3\textsuperscript{m}\ensuremath{-}3s(1)\ensuremath{-}3s(2)=3\textsuperscript{m}\ensuremath{-}3\ensuremath{\cdot}2\textsuperscript{m}+3. The guess gives the same alternating expression. These cases support, but do not prove, s(k)=\ensuremath{\Sigma}\textsubscript{j=1}\textsuperscript{k} C(k,j)(\ensuremath{-}1)\textsuperscript{k\ensuremath{-}j}j\textsuperscript{m}.
\item \begin{equation*}\mathrm{s}(1)=1;\, \mathrm{s}(2)=2^{\mathrm{m}}-2;\, \mathrm{s}(3)=3^{\mathrm{m}}-3\cdot 2^{\mathrm{m}}+3\end{equation*}
\end{enumerate}
\end{workedbox}

\end{keepsection}
\begin{keepsection}\subsection{Challenge 4.A: Surjective maps}

\begin{workedbox}{4.A Step 3a \ensuremath{\cdot} Reverse the order of a triangular sum}
\Needspace{9\baselineskip}
\begin{enumerate}
\item The left-hand iterated sum first chooses j=1,…,k, then i=1,…,j. Its index pairs are exactly \{(i,j):1\ensuremath{\leq}i\ensuremath{\leq}j\ensuremath{\leq}k\}.
\item If instead i is fixed first, i ranges from 1 to k and j ranges from i to k. Thus \ensuremath{\Sigma}\textsubscript{j=1}\textsuperscript{k}\ensuremath{\Sigma}\textsubscript{i=1}\textsuperscript{j}(…) = \ensuremath{\Sigma}\textsubscript{i=1}\textsuperscript{k}\ensuremath{\Sigma}\textsubscript{j=i}\textsuperscript{k}(…).
\item Both sides add the same term once for every pair in that finite triangular region.
\item \begin{equation*}\sum_{\mathrm{j}=1}^{\mathrm{k}}\sum_{\mathrm{i}=1}^{\mathrm{j}}(…)\, =\, \sum_{\mathrm{i}=1}^{\mathrm{k}}\sum_{\mathrm{j}=\mathrm{i}}^{\mathrm{k}}(…)\end{equation*}
\end{enumerate}
\end{workedbox}

\end{keepsection}
\begin{keepsection}\subsection{Challenge 4.A: Surjective maps}

\begin{workedbox}{4.A Step 3b \ensuremath{\cdot} Prove the binomial coefficient identity}
\Needspace{7\baselineskip}
\begin{enumerate}
\item For 0\ensuremath{\leq}i\ensuremath{\leq}j\ensuremath{\leq}k, C(k,j)C(j,i)=k!/[j!(k\ensuremath{-}j)!]\ensuremath{\cdot}j!/[i!(j\ensuremath{-}i)!].
\item Cancel j! to get k!/[i!(j\ensuremath{-}i)!(k\ensuremath{-}j)!].
\item On the other hand, C(k,i)C(k\ensuremath{-}i,j\ensuremath{-}i)=k!/[i!(k\ensuremath{-}i)!]\ensuremath{\cdot}(k\ensuremath{-}i)!/[(j\ensuremath{-}i)!(k\ensuremath{-}j)!], the same expression.
\item \begin{equation*}\mathrm{C}(\mathrm{k},\mathrm{j})\mathrm{C}(\mathrm{j},\mathrm{i})=\mathrm{C}(\mathrm{k},\mathrm{i})\mathrm{C}(\mathrm{k}-\mathrm{i},\mathrm{j}-\mathrm{i})\end{equation*}
\end{enumerate}
\end{workedbox}

\end{keepsection}
\begin{keepsection}\subsection{Challenge 4.A: Surjective maps}

\begin{workedbox}{4.A Step 3c \ensuremath{\cdot} Establish the convolution identity}
\Needspace{10\baselineskip}
\begin{enumerate}
\item Start with \ensuremath{\Sigma}\textsubscript{j=1}\textsuperscript{k} C(k,j)g(j) and insert g(j)=\ensuremath{\Sigma}\textsubscript{i=1}\textsuperscript{j} C(j,i)(\ensuremath{-}1)\textsuperscript{j\ensuremath{-}i} i\textsuperscript{m}.
\item Use Step 3a to reverse the order, then Step 3b: C(k,j)C(j,i)=C(k,i)C(k\ensuremath{-}i,j\ensuremath{-}i).
\item The result is \ensuremath{\Sigma}\textsubscript{i=1}\textsuperscript{k} C(k,i)i\textsuperscript{m} \ensuremath{\Sigma}\textsubscript{j=i}\textsuperscript{k} C(k\ensuremath{-}i,j\ensuremath{-}i)(\ensuremath{-}1)\textsuperscript{j\ensuremath{-}i}. Put r=j\ensuremath{-}i. The inner sum is \ensuremath{\Sigma}\textsubscript{r=0}\textsuperscript{k\ensuremath{-}i} C(k\ensuremath{-}i,r)(\ensuremath{-}1)\textsuperscript{r}=(1\ensuremath{-}1)\textsuperscript{k\ensuremath{-}i}.
\item It is zero unless i=k, when it equals 1. The whole expression therefore reduces to k\textsuperscript{m}, proving k\textsuperscript{m}=\ensuremath{\Sigma}\textsubscript{j=1}\textsuperscript{k} C(k,j)g(j).
\item \begin{equation*}\mathrm{k}^{\mathrm{m}}=\sum_{\mathrm{j}=1}^{\mathrm{k}}\, \mathrm{C}(\mathrm{k},\mathrm{j})\mathrm{g}(\mathrm{j})\end{equation*}
\end{enumerate}
\end{workedbox}

\end{keepsection}
\begin{keepsection}\subsection{Challenge 4.A: Surjective maps}

\begin{workedbox}{4.A Step 4 \ensuremath{\cdot} Prove s(k)=g(k) by induction}
\Needspace{10\baselineskip}
\begin{enumerate}
\item For k=1, s(1)=1 and g(1)=C(1,1)(\ensuremath{-}1)\textsuperscript{0}1\textsuperscript{m}=1.
\item Assume s(j)=g(j) for every 1\ensuremath{\leq}jm\ensuremath{-}\ensuremath{\Sigma}\textsubscript{j=1}\textsuperscript{k\ensuremath{-}1} C(k,j)s(j).
\item The convolution identity says k\textsuperscript{m}=\ensuremath{\Sigma}\textsubscript{j=1}\textsuperscript{k} C(k,j)g(j), so isolating its final term gives g(k)=k\textsuperscript{m}\ensuremath{-}\ensuremath{\Sigma}\textsubscript{j=1}\textsuperscript{k\ensuremath{-}1} C(k,j)g(j).
\item The induction hypothesis makes these right sides equal, hence s(k)=g(k). This proves the formula for every k\ensuremath{\geq}1.
\item \begin{equation*}\mathrm{s}(\mathrm{k})=\mathrm{g}(\mathrm{k})=\sum_{\mathrm{j}=1}^{\mathrm{k}}\, \mathrm{C}(\mathrm{k},\mathrm{j})(-1)^{\mathrm{k}-\mathrm{j}}\mathrm{j}^{\mathrm{m}}\end{equation*}
\end{enumerate}
\end{workedbox}

\end{keepsection}

\vfill{\small Personal worked-solution guide based on the supplied Mathematics for Engineers I Exercise Sheet 4, rewritten in my own words; no exercise text or figures are reproduced. Not an official FAU publication; check course details against the current material. Made by \textbf{Thanate Amornjiraporn (Focus)}.}

\end{document}
